\section{Positioning the Measurements}\label{sec:comparative-interpretation}

The preceding sections quantify what Motivo's compiler does and how fast it does it.
This section situates those numbers within the broader landscape of music programming tools surveyed in
Chapter~\ref{ch:background}.

\subsection{Positioning within the Music Programming Landscape}\label{detail-subsec:positioning-dsl-ecosystem}

Representative music programming systems differ along three axes that also frame this evaluation: live performance
versus offline artifacts, composition separated from synthesis, and explicit temporal structure versus implicit
sequencing (surveyed in Section~\ref{sec:related-work}).
Live-coding languages such as Sonic Pi, Tidal, Strudel, and HMusic prioritize interactive pattern manipulation and
stream events to external synthesis backends.
Real-time audio languages including ChucK and SuperCollider expose fine-grained timing and concurrent processes at the
cost of managing signal graphs and performance-time scheduling directly.
Notation and interchange formats such as ABC and MIDI encode musical information for typesetting or playback but lack
Motivo's combination of typed variables, parameterized patterns, and imperative control flow at the source level.

The benchmark evidence in this chapter does not reproduce those qualitative comparisons quantitatively.
It is instead consistent with the engineering side of Motivo's positioning: a zero-runtime compiler can expand compact
structural rules into thousands of MIDI events within the \texttt{MAX\_EVENTS} safety cap in low-millisecond time.
Where Sonic Pi advances time through \texttt{sleep} and parallel \texttt{in\_thread} blocks during performance, and
Tidal and Strudel transform time-varying patterns at live-coding time, Motivo resolves voices and loops entirely at
compile time into a portable Standard MIDI File, the trade-off already discussed in Section~\ref{sec:motivo-position}
rather than re-derived here.
The synthetic fixtures stress that offline model along every axis it exposes, nested loops, chord flattening, parallel
voices and tracks, per-note expression evaluation, and flat expansion at the compiler cap.

Unlike raw MIDI scripting, Motivo keeps compositional intent at the level of beats, patterns, and tracks; unlike ABC
notation, the source is executable and procedurally generative.
These distinctions motivated the language design; on the tested workloads, the measurements show that pursuing them
keeps compilation comfortably within interactive latency.

\subsection{Comparison with Raw MIDI Scripting}\label{detail-subsec:comparison-with-raw-midi-scripting}

The contrast between Motivo source and raw \texttt{mido}-style MIDI scripting (from \texttt{play (C4, E4, G4):2} to
the hi-hat loop in Listing~\ref{lst:midi-scripting-comparison}) is developed in
Section~\ref{sec:motivation} because it motivates the language design rather than the benchmark evidence in
this chapter.
The benchmark measurements neither quantify that verbosity gap nor depend on it; they instead confirm that pursuing
domain-level abstractions keeps compilation in the low-millisecond range across the whole suite.
